\chapter{Stability and scalar curvature}\label{sec:stability-scalar-curvature}
In dimension three the Gauss equation rewrites the Jacobi potential in terms of scalar curvature and the intrinsic Gaussian curvature. Stability then has topological consequences. The arguments follow \cite{ChMin}, Section 24, with closedness, orientability, and the convention $R_\Sigma=2K_\Sigma$ stated explicitly.

\begin{proposition}[The traced Gauss equation]\label{prop:traced-gauss}
For a smooth two-sided hypersurface $\Sigma^n\subset M^{n+1}$,
\[
 R_M=R_\Sigma+2\operatorname{Ric}_M(\nu,\nu)+|A|^2-H^2.
\]
Here scalar curvature is twice the sum of sectional curvatures over unordered orthonormal pairs, and $H$ is the sum of principal curvatures. This is \cite{ChMin}, Lemma 24.2, pp.~61--62.
\end{proposition}
\begin{proof}
At a point diagonalize the scalar second fundamental form in an orthonormal frame $e_1,\ldots,e_n$, with principal curvatures $\kappa_i$. The Gauss equation on each tangent two-plane is
$K_\Sigma(e_i,e_j)=K_M(e_i,e_j)+\kappa_i\kappa_j$.
It follows by inserting $\nabla^M_XY=\nabla^\Sigma_XY-h(X,Y)\nu$ in the curvature definition and taking the tangential component. Sum over $i\ne j$. The extrinsic sum is $H^2-|A|^2$, and the omitted ambient two-planes containing $\nu$ contribute $2\operatorname{Ric}_M(\nu,\nu)$. Rearranging proves the identity.
\end{proof}

\begin{corollary}[The Jacobi potential on a minimal surface]\label{cor:surface-jacobi-gauss}
For a two-sided minimal surface $\Sigma^2\subset M^3$,
\[
 q=|A|^2+\operatorname{Ric}_M(\nu,\nu)
 =\frac12(R_M+|A|^2)-K_\Sigma.
\]
Consequently its stability inequality is equivalent to
\[
 \frac12\int_\Sigma(R_M+|A|^2)\phi^2
 \leq\int_\Sigma\bigl(|\nabla\phi|^2+K_\Sigma\phi^2\bigr)
\]
for every compactly supported smooth test function in the interior.
\end{corollary}
\begin{proof}
Put $H=0$ and $R_\Sigma=2K_\Sigma$ in Proposition \ref{prop:traced-gauss}, solve for $q$, and substitute in Definition \ref{def:local-stability}. The factor two relates scalar and Gaussian curvature on a surface.
\end{proof}

\begin{theorem}[Topology and area of a stable surface in positive scalar curvature]\label{thm:stable-surface-positive-scalar}
Let $\Sigma$ be a nonempty smooth closed connected orientable two-sided stable minimal surface in a Riemannian three-manifold. If $R_M>0$ along $\Sigma$, then $\Sigma$ has genus zero. If moreover $R_M\geq\sigma>0$ there, then
\[
 \operatorname{Area}(\Sigma)\leq\frac{8\pi}{\sigma}.
\]
The genus conclusion is the closed orientable form of \cite{ChMin}, Theorem 24.3, attributed there to Schoen--Yau. Both the closedness and the orientability hypotheses enter the proof.
\end{theorem}
\begin{proof}
The constant test function is admissible because $\Sigma$ is closed. Corollary \ref{cor:surface-jacobi-gauss} and Gauss--Bonnet give
\[
 0<\int_\Sigma(R_M+|A|^2)
 \leq2\int_\Sigma K_\Sigma
 =4\pi\chi(\Sigma)=8\pi(1-g).
\]
Thus the nonnegative integer $g$ is zero. The same inequality with $R_M\geq\sigma$ gives $\sigma\operatorname{Area}(\Sigma)\leq8\pi$. Gauss--Bonnet and the classification of closed orientable surfaces are the topological inputs. This argument applies to an immersion as well, using its induced metric and global normal.
\end{proof}

\begin{proposition}[Rigidity of a stable torus under nonnegative scalar curvature]\label{prop:stable-torus-rigidity}
Let $\Sigma$ be a smooth closed two-sided stable minimal torus in $M^3$, and suppose $R_M\geq0$ along $\Sigma$. Then, along $\Sigma$,
\[
 A=0,\qquad R_M=0,\qquad
 \operatorname{Ric}_M(\nu,\nu)=0,\qquad K_\Sigma=0.
\]
In particular the induced metric is flat. No splitting of an ambient neighborhood is asserted.
\end{proposition}
\begin{proof}
The constant stability test and $\chi(\Sigma)=0$ give $\int(R_M+|A|^2)\leq0$. Nonnegativity and continuity force $R_M=|A|=0$. The traced Gauss equation then gives $q=\operatorname{Ric}(\nu,\nu)=-K_\Sigma$, and Gauss--Bonnet gives $Q(1,1)=-\int q=0$. For arbitrary smooth $\psi$, stability of $1+s\psi$ for both signs of small $s$ forces $Q(1,\psi)=0$. Thus $\int q\psi=0$ for every $\psi$, so $q=0$ pointwise. Hence $K_\Sigma=0$ as well. This is the equality argument in the preceding stability calculation; the polarization step is needed to obtain pointwise Ricci and Gaussian curvature conclusions.
\end{proof}

\begin{example}[A flat stable torus]\label{ex:flat-stable-torus}
In the flat product $\mathbb T^2\times\mathbb S^1$, each slice $\mathbb T^2\times\{t\}$ is two-sided and totally geodesic, with $q=0$. Hence $Q(\phi,\phi)=\int|\nabla\phi|^2\geq0$. This realizes the zero-scalar-curvature case of Proposition \ref{prop:stable-torus-rigidity}; replacing positive scalar curvature by nonnegative scalar curvature does not force a sphere.
\end{example}

\begin{example}[Why orientability is a separate hypothesis]\label{ex:stable-projective-plane}
Give $\mathbb{RP}^2$ the metric induced from the unit round sphere and consider $M=\mathbb{RP}^2\times\mathbb S^1$ with the product metric. A slice $\mathbb{RP}^2\times\{t\}$ is a closed two-sided totally geodesic minimal surface. The product direction gives a global unit normal with $\operatorname{Ric}(\nu,\nu)=0$, so the slice is stable. The ambient scalar curvature equals $2$, but the slice is not orientable and is not a sphere. This is the possibility noted in the footnote to \cite{ChMin}, Theorem 24.3. Two-sidedness alone implies orientability only when the ambient manifold is oriented.
\end{example}

\chapter{Logarithmic cutoffs and the Bernstein theorem}\label{sec:log-cutoffs-bernstein}
The following proof of the two-dimensional Bernstein theorem separates its ingredients: stability, quadratic area growth, and cutoffs of arbitrarily small Dirichlet energy. Properness supplies compact support for cutoffs defined by ambient distance. We follow \cite{ChMin}, Corollary 18.2 and Section 25, pp.~51 and 62--64.

\begin{definition}[Exhausting cutoffs with vanishing Dirichlet energy]\label{def:zero-capacity-cutoffs}
A smooth manifold without boundary has the cutoff property used here if there are compactly supported locally Lipschitz functions $\eta_j$ such that
\[
 0\leq\eta_j\leq1,\qquad \eta_j\longrightarrow1\text{ pointwise},
 \qquad \int|\nabla\eta_j|^2\longrightarrow0.
\]
Only this explicit analytic property is needed below. It is often expressed through capacity or parabolicity, but those broader equivalences are not assumed here. The logarithmic cutoffs in \cite{ChMin}, the proof of Theorem 25.1, provide a geometric example.
\end{definition}

\begin{proposition}[A nonnegative stability potential vanishes under small-energy cutoffs]\label{prop:capacity-stability-rigidity}
Suppose a smooth manifold has the cutoff property in Definition \ref{def:zero-capacity-cutoffs}, and let $q\geq0$ be a continuous function satisfying
\[
 \int q\phi^2\leq\int|\nabla\phi|^2
 \quad\text{for every }\phi\in C_c^\infty.
\]
Then $q=0$ everywhere. In particular, a two-sided stable minimal hypersurface in Euclidean space with this cutoff property is totally geodesic.
\end{proposition}
\begin{proof}
Compactly supported locally Lipschitz tests are admissible by approximation in $W^{1,2}$ on a fixed compact neighborhood of their support; $q$ is bounded there, so the potential integral also converges. Use $\eta_j$ as a test and apply Fatou's lemma to the nonnegative functions $q\eta_j^2$. This gives $\int q\leq\liminf_j\int|\nabla\eta_j|^2=0$. Continuity yields pointwise vanishing. For a Euclidean minimal hypersurface the stability potential is $q=|A|^2$.
\end{proof}

\begin{proposition}[Quadratic area growth gives logarithmic cutoffs]\label{prop:quadratic-growth-cutoffs}
Let $F:\Sigma^2\to\mathbb R^N$ be a proper smooth immersion of a surface without boundary. Suppose for some $a$ and $C<\infty$,
\[
 \operatorname{Area}\bigl(F^{-1}(B_R(a))\bigr)\leq CR^2
 \quad(R\geq1).
\]
Area is computed in the induced metric, including multiplicity of the immersion. Then $\Sigma$ has the cutoff property of Definition \ref{def:zero-capacity-cutoffs}.
\end{proposition}
\begin{proof}
Write $\rho=|F-a|$ and take $R=2^k$, $k\geq1$. Define $\eta_R$ to be one for $\rho\leq R$, zero for $\rho\geq R^2$, and $\log(R^2/\rho)/\log R$ in between. Properness makes its support compact. Since $|\nabla_\Sigma\rho|\leq1$, almost everywhere on the annulus,
\[
 |\nabla\eta_R|^2\leq\frac{1}{\rho^2(\log R)^2}.
\]
Decompose $R<\rho<R^2$ into the $k$ dyadic annuli $2^j<\rho<2^{j+1}$, $j=k,\ldots,2k-1$. On each, the integral of $\rho^{-2}$ is at most $2^{-2j}C2^{2j+2}=4C$. Hence
$\int|\nabla\eta_R|^2\leq4Ck/(k\log2)^2\to0$.
The functions take values between zero and one and converge pointwise to one. This is the logarithmic cutoff argument in \cite{ChMin}, pp.~63--64, with the squared-gradient estimate explicit.
\end{proof}

\begin{theorem}[Stable Bernstein with quadratic area growth]\label{thm:stable-bernstein-quadratic}
Let $F:\Sigma^2\to\mathbb R^3$ be a nonempty complete connected proper smooth two-sided minimal immersion without boundary. If it is stable and its induced area satisfies the quadratic growth bound in Proposition \ref{prop:quadratic-growth-cutoffs}, then its image is an affine plane.
\end{theorem}
\begin{proof}
The logarithmic cutoffs and Proposition \ref{prop:capacity-stability-rigidity} give $A=0$. The global unit normal therefore has zero derivative and is constant on connected $\Sigma$. The image lies in an affine plane and is open in that plane because $F$ is an immersion of the same dimension. Properness makes the image closed, so it is the entire plane. This is the stability-and-growth argument used for graphs in \cite{ChMin}, Theorem 25.1. The growth and properness hypotheses are part of this version; no classification of all complete stable immersions in higher dimensions is claimed.
\end{proof}

\begin{proposition}[Uniform area growth for an entire minimal graph]\label{prop:entire-graph-area-growth}
If $u\in C^2(\mathbb R^n)$ solves the scalar minimal surface equation, its graph $\Gamma$ satisfies, for every ambient center $a$ and radius $R>0$,
\[
 \mathcal H^n(\Gamma\cap B_R(a))
 \leq(n+1)\omega_{n+1}R^n.
\]
No global gradient bound is assumed. This gives an explicit constant in \cite{ChMin}, Corollary 18.2, p.~51.
\end{proposition}
\begin{proof}
The vector field $N(x,z)=(-Du(x),1)/\sqrt{1+|Du(x)|^2}$ is globally defined, has length one, and is divergence free. It is the outward normal to the subgraph $\Omega=\{z<u(x)\}$. For almost every radius, the graph meets the sphere transversely, and the divergence theorem on the bounded region $\Omega\cap B_R(a)$ gives
\[
 0=\mathcal H^n(\Gamma\cap B_R(a))
 +\int_{\Omega\cap\partial B_R(a)}N\cdot\nu_{B_R}\,d\mathcal H^n.
\]
The absolute value of the second term is no larger than the area of the entire sphere, $(n+1)\omega_{n+1}R^n$. The divergence theorem applies to this piecewise smooth boundary; its codimension-two corner has zero boundary area. Classical graph bootstrapping supplies smoothness when using Sard's theorem for the transverse radii. Approximation by good radii increasing to an arbitrary $R$ proves the estimate for every ball. This is also the flux form of the graph calibration in Proposition \ref{prop:graph-calibration}.
\end{proof}

\begin{theorem}[Bernstein's theorem for an entire graph over the plane]\label{thm:bernstein-plane-graph}
Every $C^2$ solution $u:\mathbb R^2\to\mathbb R$ of the scalar minimal surface equation is affine.
\end{theorem}
\begin{proof}
Its graph is smooth by Proposition \ref{prop:classical-graph-bootstrap}, two-sided and stable by Corollary \ref{cor:graph-positive-jacobi}, and has quadratic area growth by Proposition \ref{prop:entire-graph-area-growth}. An entire graph is proper: the inverse image of an ambient compact set is closed and has bounded horizontal coordinates. Its induced metric is complete, since it dominates the Euclidean metric on the base and is uniformly comparable to it on each compact subset. Theorem \ref{thm:stable-bernstein-quadratic} makes the image an affine plane. A plane that is a graph over all of $\mathbb R^2$ is the graph of an affine function. This is the geometric proof of \cite{ChMin}, Theorem 25.1, pp.~62--64.
\end{proof}

\begin{example}[A stable graph need not be an entire graph]\label{ex:scherk-bernstein-scope}
The nonflat Scherk graph in Example \ref{ex:scherk-graph} is stable by its positive vertical Jacobi function. Its domain is a square, so it does not satisfy Bernstein's entire-domain hypothesis. It is also not proper as that graph in $\mathbb R^3$: along $x=y\uparrow\pi/2$ its height stays zero and its points approach a finite point outside the graph. Thus pulling back ambient distance cutoffs does not automatically produce compactly supported functions on this incomplete graph. No contradiction with Theorem \ref{thm:stable-bernstein-quadratic} follows.
\end{example}
